%!TEX root = ../main.tex

\chapter{Introduction}\label{chap:introduction}

\acp{TCG} such as \ac{MTG}, \ii{Pokémon}, and various sports card collections have evolved from simple hobbies into a global market with significant economic value. The secondary market for these collectibles is vast, with prices fluctuating based on card condition, rarity, and demand. For collectors and local game stores, inventory management is a critical task. This involves sorting thousands of cards, assessing their condition, and determining their current market value. However, as collections grow, manual sorting becomes an increasingly time-consuming, tedious, and error-prone process.

And if an individual, such as a collector, but also a store or company, wants to effectively index and sort cards, the current market situation offers relatively few solutions. Moreover, these solutions are often prohibitively expensive, proprietary, or cannot be customized to individual workflow needs.

\section{Problem Statement}
The primary motivation for this work is the lack of accessible, affordable, and transparent tools for automating card indexing. Existing consumer machines often fail to meet specific requirements, are unavailable in certain geographical regions, or lack the robustness needed to handle valuable collectibles without damaging them. 

\noindent
From a technical perspective, the challenge is twofold:
\begin{enumerate}[(a)]
    \item \ii{Mechanical:} Designing a sorting mechanism that handles delicate paper objects reliably without jamming or causing surface damage.
    \item \ii{Computational:} Developing a computer vision pipeline capable of fine-grained classification. Unlike standard object detection, trading cards often share identical layouts with only minor variation in artwork or text, requiring high-precision feature extraction and matching under varying lighting conditions.
\end{enumerate}

\section{Objectives}
The goal of this project is to develop an automated system for recognizing and classifying trading cards that follows the philosophy of Open-source software and hardware. By designing a robot based on accessible components, such as an Arduino microcontroller and stepper motors, and combining it with modern computer vision techniques, this work aims to democratize the technology needed for card sorting.

%%%%%% FROM MRS TEMPLATE:
\section{Thesis Contribution}


\section{Thesis Structure}


\section{Mathematical notation}
It is a good practice to define basic mathematical notation in the introduction.
See \reftab{tab:mathematical_notation} for an example.
\begin{table*}[ht!]
  \scriptsize
  \centering
  \noindent\rule{\textwidth}{0.5pt}
  \begin{tabular}{lll}
    $\mathbf{x}$, $\bm{\alpha}$ & vector, pseudo-vector, or tuple\\
    $\mathbf{\hat{x}}$, $\bm{\hat{\omega}}$& unit vector or unit pseudo-vector\\
    $\mathbf{\hat{e}}_1, \mathbf{\hat{e}}_2, \mathbf{\hat{e}}_3$ & elements of the \emph{standard basis} \\
    $\mathbf{X}, \bm{\Omega}$ & matrix \\
    $\mathbf{I}$ & identity matrix \\
    $x = \mathbf{a}^\intercal\mathbf{b}$ & inner product of $\mathbf{a}$, $\mathbf{b}$ $\in \mathbb{R}^3$\\
    $\mathbf{x} = \mathbf{a}\times\mathbf{b}$ & cross product of $\mathbf{a}$, $\mathbf{b}$ $\in \mathbb{R}^3$\\
    $\mathbf{x} = \mathbf{a}\circ\mathbf{b}$ & element-wise product of $\mathbf{a}$, $\mathbf{b}$ $\in \mathbb{R}^3$ \\
    $\mathbf{x}_{(n)}$ = $\mathbf{x}^\intercal\mathbf{\hat{e}}_n$ & $\mathrm{n}^{\mathrm{th}}$ vector element (row), $\mathbf{x}, \mathbf{e} \in \mathbb{R}^3$\\
    $\mathbf{X}_{(a,b)}$ & matrix element, (row, column)\\
    $x_{d}$ & $x_d$ is \emph{desired}, a reference\\
    $\dot{x}, \ddot{x}, \dot{\ddot{x}}$, $\ddot{\ddot{x}}$ & ${1^{\mathrm{st}}}$, ${2^{\mathrm{nd}}}$, ${3^{\mathrm{rd}}}$, and ${4^{\mathrm{th}}}$ time derivative of $x$\\
    $x_{[n]}$ & $x$ at the sample $n$ \\
    $\mathbf{A}, \mathbf{B}, \mathbf{x}$ & LTI system matrix, input matrix and input vector\\
    \emph{SO(3)} & 3D special orthogonal group of rotations\\
    \emph{SE(3)} & \emph{SO(3)}~$\times~\mathbb{R}^3$, special Euclidean group\\
  \end{tabular}
  \noindent\rule{\textwidth}{0.5pt}
  \caption{Mathematical notation, nomenclature and notable symbols.}
  \label{tab:mathematical_notation}
\end{table*}
